\centerline{\bf \Huge LTE, 4G, H+...}

\begin{flushright}
        \scriptsize{Эльзята тýпик\\(с) Какашка}
\end{flushright}

\problem{1}
\textbf{LTE-лемма в случае $p=2$.}
Пусть $x$ и $y$ - различные нечетные целье числа и $4 \mid x-y$. Тогда для любого натурального $n$ выполнено

$$
\left\|x^n-y^n\right\|_2=\|x-y\|_2+\|n\|_2 .
$$

\problem{2}
С помощью LTE-леммы узнайте на сколько нулей заканчивается число $4^{5^6}+6^{5^4}$?

\problem{3}
Пусть натуральные числа $x, y, p, n, k$ таковы, что $x^n+y^n=p^k$. Докажите, что если число $n>1$ - нечетное, а число $p$ - простое нечетное, то $n$ является степенью числа $p$ с натуральным показателем.

\problem{4}
Докажите, что для любого натурального $a>2$ найдется такое натуральное $n>1$, что $a^n-1$ делится на $n^2$. Верно ли это утверждение для $a=2$ ?

\problem{5}
На доске написаны $2 n$ последовательных целых чисел. За ход можно разбить написанные числа на пары произвольным образом и каждую пару чисел заменить на их сумму и разность (не обязательно вычитать из большего числа меньшее, все замены происходят одновременно). Докажите, что на доске больше никогда не появятся $2 n$ последовательных целых чисел.
%ММО2020-11.1.6

\problem{6}
Для натурального $n$ обозначим $S_n=1!+2!+\ldots+n!$. Докажите, что при некотором $n$ у числа $S_n$ есть простой делитель, больший $10^{2025}$.
%ВСОШ2012-11.8

\problem{7}
Найдите все пары натуральных чисел $(m,n)$ таких, что 
$$
    \left(2^n-1\right)\left(2^n-2\right)\left(2^n-4\right) \cdots\left(2^n-2^{n-1}\right)=m!
$$
%IMOSL2019-N1

\problem{8}
Найдите все тройки натуральных чисел $(a, b, p)$, где $p$ простое, такие что:

$$
    a^p=b!+p .
$$
%IMOSL2022-N4